% !TeX root = ../papers/nc-p3.tex
\section{相対trace格子の核の冪零性と一つの次元ベクトルでの検査}
\label{sec:relative-trace-nilpotence}
\begin{theorem}\label{thm:relative-trace-nilpotence}
前節の正則特殊化で $B=\mathcal B/t\mathcal B$、$T=\mathcal B C$ とすると、
$B\to T/tT$ の核は冪零両側イデアルである。従って整域でない正則点には経路長4または8の
非零な\emph{冪零中心元}がある。各 $j=4,8$ で、次元ベクトル $(4,4,4,4)$ の関係付き四窓行列の
被約な表現多様体上に、零番でのtraceによる線型写像
\[
         Z(B)_j\longrightarrow k[\operatorname{Rep}_{4\delta}(B)_{\rm red}]
\]
を取る。十五次元閉包の有限四辺の正則点では、両写像の単射性と合成整域性は同値である。
従って両写像が単射な有限Frobenius候補では84次検査が全次数正則性自体を判定する。
ここで使う二つのtrace判定は非斉次表現の有限行列式と可換な根基イデアルで完結する。
\end{theorem}
\begin{proof}[一つの非斉次右加群による核の検出]
$K=\ker(B\to T/tT)$ とし、その正次数の任意の斉次元を $a$、斉次な持ち上げを
$\tilde a\in\mathcal B\cap tT$ とする。零次数には頂点冪等元の独立性により核はない。
右加群 $M=\mathcal B/(1-\tilde a)\mathcal B$ を作る。$1-a$ の左乗法は正次数部分と最低次数を比較して
$B$ 上単射である。従って $M$ は $t$ 平坦である。実際 $tb=(1-\tilde a)d$ なら剰余で $\bar d=0$ となり
$t$ で割れる。$t$ 進完備化 $\widehat M$ も $t$ 平坦で、$\widehat M/t\widehat M=M/tM$ である。
各厚みでの $t$ 倍の核・像を取り、逆極限で一つずつ割ればよい。

$I=\{q\in\mathcal B:Tq\subset\mathcal B\}$ を取る。これは $\mathcal B$ の両側イデアルである。
一般点で全行列環の次元 $16^2$ 個を交代にしたCapelli多項式の全ての $\mathcal B$ 値は $I$ に入る。
実際trace挿入を反復すると各値とtraceの任意の積が $\mathcal B$ に入り、$T=C\mathcal B$ である。
$M$ の剰余生成元を $v$ とすると $v=v\tilde a^m$。各 $b\in\mathcal B$、$q\in I$ について
$v bq=v\tilde a^m bq\in t^m M$ である。$t^{-m}\tilde a^m\in T$、$bq\in I$ を使う。
従って $MI$ は完備化で零になる。

$\mathcal B_\eta$ は既知の一般中心上有限で、Azumaya部分の各閉点ではこの交代値イデアルは全体になる。
実際その行列繊維で単位行列環を張り、全行列単位の交代評価は非零で、生成する両側イデアルは全体である。
一般の $I_\eta$ は分母を払えるのでこれらの全ての評価を含む。よって $\mathcal B_\eta/I_\eta$ の台は
既出の三十の特異射影点に対応する錐直線と頂点に含まれる。さらに右加群
\[
        (\mathcal B_\eta/I_\eta)/(1-\tilde a)(\mathcal B_\eta/I_\eta)
\]
の台はこの集合の零次元部分だけである。各錐直線の非頂点部分は次数作用の一軌道である。
一指標の有限次元左正則行列における $1-\lambda^{\deg a}\tilde a$ の行列式は
$\lambda$ の多項式で定数項1を持ち、可逆性が破れる倍率は有限個しかない。これは基礎体の代数閉拡大で調べてよい。
各繊維で単位になればその指標は台にない。従って表示した加群は一般の底体上有限次元である。
その像は $M$ の $\widehat M[1/t]$ 内の像を張る。特殊 $M/tM$ の独立な元の持上げは
この一般空間でも独立である。係数の最小付値を除いて $\widehat M$ で還元すればよい。
これで $M/tM=B/(1-a)B$ が有限次元と分かる。

この有限次元加群の射影次元は $1-a$ の単射により高々1である。しかし非零の有限次元加群 $D$ は
四次元Calabi--Yau双対により $\operatorname{Ext}^4_B(D,D)=\operatorname{End}_B(D)^*\ne0$ を満たす。
よって $B/(1-a)B=0$。$1-a$ の右逆を最低次数から比較すると有限和 $1+a+\cdots$ となるので $a$ は冪零である。
ここまで $a$ の中心性を使っていない。任意の有限次元単純表現の上で、$K$ の各斉次元のtraceは零となる。
$K$ の像は行列環の両側イデアルなので非零なら全体となり、全trace零とは両立しない。
アフィンPI代数の全単純は有限次元だから $K$ はJacobson根基に入り、Braunの冪零性により $K$ 自体も冪零となる。
\end{proof}

\begin{proof}[二つの短い中心次数]
前節の $c=\overline{t^u\tau(b)}\ne0$ は $u>0$ かつ $\tau(b)\in T$ なので $K$ に入り、
中心冪零元を得る。全ての行列表現で各頂点成分も冪零行列に作用し、被約な表現多様体上でtraceは恒等的に零である。
逆に整域の正則点は中心上有限であり、一般Azumaya部分には各頂点次元4の単純がある。
任意の非零中心元はその稠密開集合のある閉点で非零なスカラーに作用するため、零番traceが消えない。
これで同値を得る。実装上は、四辺それぞれの四矢に四行四列の未知行列を置き、二次関係の行列成分の生成イデアルの
可換根基を取る。中心とそのtraceの計算は経路長8と、矢との可換性検査のための9までで足りる。
各次数の中心基底のtrace多項式がこの根基を法として線型独立かを判定すればよい。
\end{proof}
